% Compiler: pdfLaTeX
\documentclass[11pt,a4paper]{article}

\usepackage[utf8]{inputenc}
\usepackage[T1]{fontenc}
\usepackage{lmodern}
\usepackage{microtype}
\usepackage{geometry}
\usepackage{enumitem}
\usepackage{titlesec}
\usepackage{fancyhdr}
\usepackage{hyperref}

% Document details
\newcommand{\researchtitle}{Survivorship bias}
\newcommand{\researchshorttitle}{Survivorship bias}
\newcommand{\researchdate}{2026-10-25}

% Layout
\geometry{left=2.7cm,right=2.7cm,top=2.5cm,bottom=2.5cm,headheight=14pt}
\hypersetup{hidelinks,pdftitle={\researchtitle}}
\urlstyle{same}
\frenchspacing
\setlength{\parindent}{0pt}
\setlength{\parskip}{0.55em}
\linespread{1.04}
\setlength{\emergencystretch}{1em}

\titleformat{\section}
    {\large\bfseries}{\thesection}{0.65em}{}
\titleformat{\subsection}
    {\normalsize\bfseries}{\thesubsection}{0.65em}{}
\titlespacing*{\section}{0pt}{1.15em}{0.35em}
\titlespacing*{\subsection}{0pt}{0.7em}{0.2em}
\setlist[itemize]{leftmargin=1.35em,itemsep=0.25em,topsep=0.2em}
\setlist[enumerate]{leftmargin=1.6em,itemsep=0.25em,topsep=0.2em}

\pagestyle{fancy}
\fancyhf{}
\fancyhead[L]{\small Random Research}
\fancyhead[R]{\small\researchshorttitle}
\fancyfoot[C]{\small\thepage}
\renewcommand{\headrulewidth}{0pt}
\renewcommand{\footrulewidth}{0pt}

\begin{document}

{\LARGE\bfseries\researchtitle\par}
\vspace{0.4em}
{\small\researchdate\par}

\section{Overview}

Survivorship bias appears when the cases still visible at the end are a
selective part of the group that existed at the beginning. Looking only
at those cases can give a misleading account of what usually happens.
\cite{lane}

Success stories are an easy place to notice the problem. We can inspect
the winner's habits, while people who followed similar habits and failed
may have disappeared from view.

\section{Findings}

Lane's statistics text uses investment funds to explain the mechanism.
Poorly performing funds can close or merge. Calculating historical
performance using only funds still operating leaves some of those
earlier options out of the sample. \cite{lane}

This is a problem that research methods can address. The SPIVA methodology
includes funds that liquidated or merged, since they belonged to the
choices available at the beginning of the period. Its comparison starts
with that original group. \cite{spiva}

The issue is how cases enter or leave the evidence. A large collection
of surviving cases can still omit the group needed to answer the
question. \cite{lane}

\section{Example}

Imagine a writing course advertises interviews with ten former students
who became successful authors. Their accounts might contain useful advice.
But the interviews cannot, on their own, show how likely a new student is
to become successful.

For that question, I would also want to know how many people took the
course, what happened to the others, and what counts as success.
This is a hypothetical example. The missing information prevents a fair
estimate even if every interview is truthful.

\section{Discussion}

I think survivorship bias is easy to overlook because the available
stories often sound convincing. They include specific decisions and
people we can identify with. The missing cases rarely provide an equally
memorable story.

I would still read an experienced person's account. It can suggest ideas
worth trying. I would be more cautious when the account becomes a claim
that following those steps reliably produces the same outcome.

There is also a limit to the correction: failure stories alone are
selective too. Replacing a set of winners with a set of failures does
not establish what happened to the whole group.

\section{Practical implications}

Before drawing a lesson from a result, I would ask who was eligible to
appear in the evidence and who disappeared before it was collected.
Whenever possible, I would look for a group defined at the beginning
and followed over time.

If that information is unavailable, I would treat the story as a source
of possibilities and leave the success rate unresolved.

\renewcommand{\refname}{References}
\begin{thebibliography}{9}
\small
\setlength{\itemsep}{0.35em}

\bibitem{lane}
Lane, D. M.
\href{https://onlinestatbook.com/2/research_design/sampling.html}
{Sampling bias.}
\textit{Online Statistics Education}.

\bibitem{spiva}
S\&P Dow Jones Indices (2025).
\href{https://www.spglobal.com/spdji/en/index-tv/article/what-is-spiva/}
{What is SPIVA?}
Methodology overview and transcript, September 10.

\end{thebibliography}
\end{document}
